%------------------------------------------------------------------------------%
\begin{exercises}

%------------------------------------------------------------------------------%
\begin{exercise}
Usando frações parcias, calcule as seguintes integrais.
\begin{mathlist}
  \item \int \frac{x-9}{(x+5)(x-2)} dx
  \item \int_0^1 \frac{2}{2x^2+3x+1} dx
  \item \int_3^4 \frac{x^3-2x^2-4}{x^3+2x^2} dx
  \item \int \frac{x^2+1}{(x-3)(x-2)^2} dx
  \item \int \frac{10}{(x-2)(x^2+9)} dx
  \item \int \frac{x^3-x+6}{x^3+3x}dx
  \item \int  \frac{x^2-2x-1}{(x-1)^2(x^2+1)}dx
  \item \int \frac{1}{x^3-1}dx
  \item \int \frac{1}{x(x^2+4)^2}dx
  \item \int \frac{x^2+2x+1}{x^3-x}dx
 %\item \int \sin(7x) \cos(5x)dx
 %\item \int\frac{\sin(\phi)}{\cos^3(\phi)} d \phi
 %\item \int \tan^5(x)dx
 %\item \int_{\frac{\pi}{6}}^{\frac{\pi}{2}} \cot^2(x)dx
\end{mathlist}
\begin{answer}
\begin{mathlist}
  \item 2 \ln \vert x+5 \vert-\ln \vert x-2 \vert+C
  \item 2 \ln(\frac{3}{2})
  \item \frac{7}{6}+ \ln(\frac{2}{3})
  \item 10\ln \vert x-3 \vert-9 \ln \vert x-2 \vert+ \frac{5}{x-2}+C
  \item \frac{3}{8} \pi
  \item 2\ln \vert x \vert-\frac{1}{2}\ln(x^2+3)
        -\frac{1}{\sqrt{3}}\arctan(\frac{x}{\sqrt{3}})+C
  \item \ln \vert x-1 \vert+ \frac{1}{x-1}-\frac{1}{2} \ln(x^2+1)+ \arctan(x)+ C
  \item \frac{1}{3}\ln \vert x-1\vert- \frac{1}{2}\ln(x^2+x+1)-\frac{1}{\sqrt{3}}
        \arctan \left(\frac{1}{\sqrt{3}} (2x+1) \right)+C
  \item \frac{1}{16} \ln \vert x\vert-\frac{1}{32}\ln(x^2+4)+\frac{1}{8(x^2+4)}+C
  \item \ln \vert x \vert -\ln \vert x-1 \vert+ \ln \vert x-1\vert +C
\end{mathlist}
\end{answer}
\end{exercise}

%------------------------------------------------------------------------------%
\begin{exercise}
Faça uma substituição simples para expressar o integrando como uma função
racional e então calcule a integral \(\int \frac{e^{2x}}{e^{2x}+3e^x+2}dx\).
\begin{answer}
  $\ln \left( \frac{(e^x+2)^2}{e^x+1} \right)+ C$
\end{answer}
\end{exercise}

%------------------------------------------------------------------------------%
\begin{exercise}
Faça uma substituição simples para expressar o integrando como uma função
racional e então calcule a integral \(\int \frac{e^{x}}{(e^x-2)(e^{2x}+1)}dx\).
\end{exercise}

%------------------------------------------------------------------------------%
\begin{exercise}
  Use integração por partes, juntamente com frações parciais para resolver a
  integral $\int \ln(x^2-x+2)dx$.
\begin{answer}
  $(x-\frac{1}{2})\ln(x^2-x+2)-2x+\sqrt{7} \arctan
  \left(\frac{2x-1}{\sqrt{7}}\right)+C$
\end{answer}
\end{exercise}

%------------------------------------------------------------------------------%
\begin{exercise}
  Encontre a área da região sob a curva
  $y=\frac{x^2+1}{3x-x^2}$ de $x=1$ até $x=2$.
\begin{answer}
$ A = -1 + \frac{11}{3} \ln 2 $
\end{answer}
\end{exercise}

\begin{exercise}
Determine se cada integral é convergente ou divergente.
Calcule aquelas que são convergentes.
\begin{mathlist}
  \item \int_{0}^{\infty} \frac{1}{v^2+2v-3} dv
  \item \int_{0}^{3} \frac{dx}{x^2-6x+5} dx
\end{mathlist}
\begin{answer}
\begin{alphalist}
  \item converge para \(\frac{1}{4} \ln 5\)
  \item Divergente
\end{alphalist}
\end{answer}
\end{exercise}

\begin{exercise}A integral \(\int \frac{\arctan(x)}{x^2}dx\) pode ser resolvida usando
  primeiro uma integração por partes, depois frações parciais e por fim, é
  necessário usar uma substituição simples. Resolva essa integral. Além disso, mostre que a integral $\int_1^\infty \frac{\arctan(x)}{x^2}dx$ é convergente e vale $\frac{\pi}{4}+ \frac{1}{2}\ln2$.
\end{exercise}
\emph{Revisão de Técnicas de Integração}

Os dois próximos exercícios são muito importantes para que vocês possam pegar a
manha de qual técnica utilizar, uma vez que você tera que decidir qual(is)
técnicas usar.

%------------------------------------------------------------------------------%
\begin{exercise}
Nos itens a seguir, decida qual técnica usar e depois calcule a integral.
\begin{mathlist}
  \item \int_{0}^1 \frac{e^{\arctan(x)}}{1+x^2} dx
  \item \int \sin^5(x) \cos^4(x) dx
  \item \int \frac{dx}{(1-x^2)^{\frac{3}{2}}} dx
  \item \int \frac{\cos^3(x) \sin(x)}{(9-\cos^2(x))^{\frac{5}{2}}} dx
  \item \int_1^3 r^4 \ln(r) dr
  \item \int_0^4 \frac{x-1}{x^2-4x+5}dx
  \item \int  \frac{x^2}{\sqrt{1+x^2}}dx
  \item \int \sqrt{3-2x-x^2} dx
  \item \int_0^{\frac{\pi}{4}} \tan^3(x)\sec^2(x) dx
  \item \int \arccos(x) dx
  \item \int t \sin(t) \cos(t)dt
  \item \int\frac{dx}{x^2 \sqrt{4x^2-1}}dx
  \item \int \tan^5(x)dx
  %\item \int_{\frac{\pi}{6}}^{\frac{\pi}{2}} \cot^2(x)dx$
\end{mathlist}
\begin{answer}
  \begin{mathlist}
    \item e^{\frac{\pi}{4}}-1
    \item -\frac{1}{5} \cos^5(x)+\frac{2}{7}\cos^7(x)-\frac{1}{9} \cos^9(x)+C
    \item \frac{x}{\sqrt{1-x^2}}+C
    \item \;
    \item \frac{243}{5}ln3-\frac{243}{25}
    \item -\frac{1}{3}ln5
    \item \frac{1}{2}(1+x^2)^{\frac{3}{2}}-(1+x^2)^{\frac{1}{2}}+ C
    \item 2\arcsin(\frac{x+1}{2})+ \frac{x+1}{2}\sqrt{3-2x-x^2} +C
    \item \frac{1}{4}
    \item x\arccos(x)-\sqrt{1-x^2} +C
    \item -\frac{1}{4}t \cos(2t)+ \frac{1}{8} \sin(2t)+C
  \end{mathlist}
\end{answer}
\end{exercise}

%------------------------------------------------------------------------------%
\begin{exercise}
Calcule as Integrais. Em cada um dos itens você poderá de usar mais de uma
técnica ou até a mesma técnica mais de uma vez,
\begin{mathlist}
  \item\int_0^{\pi} t \cos^2(t) dt
  \item\int x \arccos(x)dx
  \item\int \frac{cos^3(\sqrt{x})}{\sqrt{x} \hspace{0,1cm} \csc(\sqrt{x})}dx
  \item\int \frac{r}{\sqrt{r^2-2r}} dr
  \item\int x^5e^{-x^3} dx
  \item\int \frac{dx}{e^{2x}-3e^x}dx
  \item\int \frac{\sec(\theta)\tan(\theta)}{\sec^2(\theta)- \sec(\theta)}d \theta
  \item\int \frac{\ln(x)}{x \sqrt{1+(\ln(x))^2}}dx
  \item\int \frac{e^x}{(e^{2x}+6e^x+10)^{\frac{3}{2}}} dx
  \item\int \frac{3x^2-2}{x^3-2x-8}dx
\end{mathlist}
\end{exercise}

%------------------------------------------------------------------------------%

%------------------------------------------------------------------------------%
%\begin{exercise}
%\begin{answer}
%\end{answer}
%\end{exercise}

%------------------------------------------------------------------------------%
%\begin{exercise}
%\begin{mathlist}[2]
%  \item
%  \item
%  \item
%\end{mathlist}
%\begin{answer}
%  \begin{alphalist}[3]
%    \item
%    \item
%    \item
%  \end{alphalist}
%\end{answer}
%\end{exercise}

\end{exercises}
%------------------------------------------------------------------------------%
